\section{Study design}
\label{sec:methods}
The study developed in stages. We first examined two ways of reducing repeated
reasoning after changes to a knowledge base, then used the observed costs and
the literature review to select further experiments. The later questions and
comparisons were therefore exploratory. Protocols were recorded before their
respective timing series, but the research programme was not publicly
preregistered. Earlier trials, unsuccessful cases, and subsequent replications
are retained in the supplementary evidence.

\subsection{Interventions and comparisons}
The first intervention changes how groups of records satisfying a rule's
conditions are combined. One option first collects the records affected by a
change; another starts with records meeting all the conditions and then checks
which were affected. We investigated whether information about recent changes
helps choose between these equivalent approaches, including the cost of the choice.

The second intervention shares explanations for conclusions that remain valid
after a change. Several individuals may satisfy exactly the same relevant
conditions. When they do, one calculation can establish which conclusions
still follow for all of them. These explanations can prevent unnecessary
withdrawal and restoration of conclusions. If no suitable explanation is
found, the ordinary maintenance procedure still determines the answer. This
restricted use of shared explanations builds on support-aware maintenance and
shared reasoning~\cite{motik2015bf,hu2018,hu2019compressed,hu2022modular}.

Follow-up changes reduced repeated work in three further ways: reusing checks
that remain valid when information is added, updating access structures only
where information has changed, and preparing common rule operations for
repeated use. These changes apply established techniques; the experiment
evaluates their combined effect. A later experiment reconsidered evaluation
order by briefly examining likely consequences of alternative choices. This
was motivated by recent work on avoiding large intermediate results
\cite{qiao2026,bekkers2025}. The examination changes the search order, without
excluding possible answers.

We distinguish three comparisons. Controlled variants of the same system
estimate an intervention's effect. Comparisons with a preserved earlier
version describe cumulative progress. Executing an independent reasoner on
matched inputs and tasks tests relative performance beyond the local system.
Published timings on different machines provide context, but cannot supply
speed ratios for our experiments.

\subsection{Controlled and established workloads}
The initial experiment combined the first two interventions in a factorial
design: each was either present or absent, yielding four variants. Eighteen
constructed cases varied the concentration of changes, surviving explanations,
recursion, and opportunities to share reasoning. Nine repetition blocks tested
all variants in fresh processes, with a rotating order. Complete update time
was the primary outcome. A separate replication used new random seeds and
144 fresh processes to re-examine the main explanation-sharing benefit and its
adverse cases. It supplements the original observations rather than replacing
them.

External validation began with LUBM, which
provides a university ontology, generated data, and reference answers
\cite{guo2005,lubmwebsite}. We used one generated university and retained the
logical ontology. Three repetition blocks compared three local variants, and
all fourteen queries were checked against the complete official answer sets.
This tests established questions on a modest dataset; it is not a scalability
study.

For changing rules, we used the positive university rule programme supplied
with ZodiacEdge~\cite{xu2023,zodiacartifact}. Two scenarios added and removed
designated sets of eight and sixteen rules. Four repetition blocks compared
four local variants; a separate four-block comparison alternated our candidate
and ZodiacEdge itself. Additional samples removed and restored 1,000 facts,
following the update size used in earlier maintenance research~\cite{hu2018}.
Samples within one process are repeated observations, not independent process
replicates. The author's rule programme and the complete LUBM ontology are
distinct reasoning tasks.

The Transaction Processing Performance Council's Benchmark H (TPC-H) evaluates
database systems answering analytical questions over business data~\cite{tpch}.
Our query experiments used the conditions underlying its queries Q3, Q5,
and Q9 at two data scales. These six components retain the combinations of
input rows satisfying each query, but exclude aggregation and result ordering.
They therefore do not constitute a full TPC-H evaluation or reproduce the
larger study of Qiao and colleagues~\cite{qiao2026}. The first comparison used
three blocks, with all three components evaluated sequentially in each fresh
process. A follow-up gave each component its own process and compared three
variants across three blocks. It added controls deliberately chosen to expose
unhelpful decision costs and misleading intermediate patterns, giving 90
processes overall. The follow-up revisits the workload that motivated it;
it does not provide independent evidence of generalization.

\subsection{Expressiveness, correctness, and measurement}
A separate experiment compared our reasoner with HermiT~\cite{glimm2014hermit}
on a finite family using existential statements: every member of one class
must relate to some member of another. Both systems determined consistency
and returned all named members of two target classes. Populations of 100,
1,000, and 10,000 individuals were tested in three blocks, comparing the
candidate, its earlier baseline, and HermiT. These 27 processes test a selected
DLP L3 task beyond the structural OWL 2 RL syntax. The family also belongs to
OWL 2 EL~\cite{owlprofiles2012}, and its answers do not require explicitly
constructing the unnamed individuals. It tests computational specialization,
without assuming that such ontologies are common.

Correctness was checked separately from performance. Update experiments
compared complete conclusions with fresh recomputation; small cases also used
independently written reference procedures. Query experiments checked exact
answers against official LUBM answers or independent database evaluation.
The existential experiment used independently justified answers and controls
omitting each essential statement. Agreement of counts alone was insufficient.

A further comparison returned to the original Bach examples at their original
scale~\cite[Table~2.5, p.~35; Figure~6.2, p.~150]{volz2004}. The current
reasoner, its preserved earlier version, and HermiT answered 25 questions on
the complete ontology. A separate family task asked nine questions, including
the complete ancestor and dynasty relationships. All four systems evaluated
the original and three changed family states from scratch; the two local
versions and ZodiacEdge also applied and reversed each change while retaining
state. Twelve repetition blocks specified 336 fresh processes, with rotating
system order. The older version was measured anew: its saved historical
results contained no Bach comparison. These small examples test source fidelity
and update behaviour, without making a scaling claim.

For this comparison, the complete ontology's existence requirements and
constraints fall outside ZodiacEdge's positive-rule task. Its family evaluation
uses the same recorded relationships and rules. A wrong answer disqualifies
the entire system--task combination from a timing ranking; passing executions
are not selected to hide unsuccessful ones. Separate diagnostic checks
investigate discrepancies while retaining the recorded implementation.

Timings include complete updates or complete consumption of query answers.
Query computation and data preparation are reported separately, with combined
costs where measured. HermiT comparisons include loading, consistency, and
answer retrieval because systems perform inference at different stages.
Independent checks lie outside these task intervals. In controlled comparisons,
ratios pair observations within repetition blocks and report their median,
rather than dividing two separately calculated medians. Historical comparisons
use ratios of the medians from separate runs.

All new measurements used one Apple M4 Max desktop. Process restarts did not clear
filesystem caches. Unrelated activity may have overlapped the primary study;
its extent is unknown, and the original samples were retained. Later timed
studies were scheduled separately. Small repetition counts, one host, and
selected workloads make the comparisons descriptive. Detailed protocols,
timing boundaries, individual observations, and preservation arguments are
provided in the supplement.
